\documentclass[conference]{IEEEtran}
\usepackage{cite}
\usepackage{amsmath,amssymb,amsfonts}
\usepackage{algorithmic}
\usepackage{graphicx}
\usepackage{textcomp}
\usepackage{xcolor}
\usepackage{hyperref}
\usepackage{textgreek}
\usepackage[backend=biber, style=ieee]{biblatex}

\hypersetup{
    colorlinks=true,
    linkcolor=blue,
    citecolor=green,
    urlcolor=cyan,
}

\def\BibTeX{{\rm B\kern-.05em{\sc i\kern-.025em b}\kern-.08em
    T\kern-.1667em\lower.7ex\hbox{E}\kern-.125emX}}

\title{{{TITLE}}}
\author{\IEEEauthorblockN{{{AUTHOR}}}}

\begin{document}

\maketitle

\begin{abstract}
%% -- BLOCK: abstract -- %%
This is a placeholder for the abstract. The abstract should be a concise summary of the paper's contributions.
%% -- ENDBLOCK: abstract -- %%
\end{abstract}

\begin{IEEEkeywords}
%% -- BLOCK: keywords -- %%
Enter key terms or phrases in alphabetical order, separated by commas. For a list of suggested keywords, send a blank e-mail to keywords@ieee.org or visit \url{http://www.ieee.org/organizations/pubs/ani_prod/keywrd98.txt}
%% -- ENDBLOCK: keywords -- %%
\end{IEEEkeywords}

\section{Introduction}
%% -- BLOCK: introduction -- %%
This is a placeholder for the introduction. This section should provide a clear background, state the problem, and outline the contributions of the paper.
%% -- ENDBLOCK: introduction -- %%

\section{Related Work}
%% -- BLOCK: related_work -- %%
This is a placeholder for the related work. Discuss existing literature and how this work differs or builds upon it.
%% -- ENDBLOCK: related_work -- %%

\section{Methodology}
%% -- BLOCK: methodology -- %%
This is a placeholder for the methodology. Describe the methods, algorithms, and experimental setup.

Example of a figure:
\begin{figure}[htbp]
\centerline{\includegraphics[width=\columnwidth]{figure.png}}
\caption{Example of a figure caption.}
\label{fig}
\end{figure}
%% -- ENDBLOCK: methodology -- %%

\section{Results and Discussion}
%% -- BLOCK: results_discussion -- %%
This is a placeholder for the results and discussion. Present the results and interpret them.

Example of a table:
\begin{table}[htbp]
\caption{Table Type Styles}
\begin{center}
\begin{tabular}{|c|c|c|c|}
\hline
\textbf{Table}&\multicolumn{3}{|c|}{\textbf{Table Column Head}} \\
\cline{2-4} 
\textbf{Head} & \textbf{\textit{Table column subhead}}& \textbf{\textit{Subhead}}& \textbf{\textit{Subhead}} \\
\hline
copy& More table copy$^{\mathrm{a}}$& &  \\
\hline
\multicolumn{4}{l}{$^{\mathrm{a}}$Sample of a Table footnote.}
\end{tabular}
\label{tab1}
\end{center}
\end{table}
%% -- ENDBLOCK: results_discussion -- %%

\section{Conclusion}
%% -- BLOCK: conclusion -- %%
This is a placeholder for the conclusion. Summarize the paper's findings and suggest future work.
%% -- ENDBLOCK: conclusion -- %%

\section*{Acknowledgment}
%% -- BLOCK: acknowledgment -- %%
This is a placeholder for the acknowledgment. Acknowledge funding sources or other contributors here.
%% -- ENDBLOCK: acknowledgment -- %%

%% -- BLOCK: bibliography -- %%
% This block is for BibTeX entries.
% Example:
% @article{einstein,
%   author = {Albert Einstein},
%   title = {Zur Elektrodynamik bewegter Körper},
%   journal = {Annalen der Physik},
%   volume = {322},
%   number = {10},
%   pages = {891--921},
%   year = {1905},
%   doi = {10.1002/andp.19053221004}
% }
%% -- ENDBLOCK: bibliography -- %%

\section*{References}
\addcontentsline{toc}{section}{References}
%% -- BLOCK: references -- %%
\printbibliography
%% -- ENDBLOCK: references -- %%

\end{document}
